\section{Framework-Defined Infrastructure：把应用当成高层语言}

\term{Framework-Defined Infrastructure}（FDI）把框架约定视为高层语言。框架知道哪些 route 是静态的、哪些需要函数、哪些数据可以缓存、哪些请求需要 middleware；平台在 build time 解析这些语义，生成底层基础设施配置。与 Infrastructure as Code 不同，开发者不直接声明每个云资源，而是声明应用，基础设施成为编译输出。

\subsection{Cloud compiler 与 intermediate representation}

设应用源码为 $A$，框架语义抽取器为 $S_F$，平台编译器为 $C_P$，则部署计划可抽象为
\[
I=C_P\bigl(S_F(A)\bigr),
\]
其中 $S_F(A)$ 不是原始源码，而是 \term{intermediate representation}（IR，中间表示）：route graph、static assets、function boundaries、cache policy、runtime hints 和环境元数据。IR 让平台不必理解每个框架内部实现，只需要稳定的输出契约。

\lecturefigure{03-cloud-compiler.png}{Framework semantics 经过平台 IR 转换为部署计划。}{本地字幕 04:40--06:00；Vercel 官方 Framework-Defined Infrastructure。}

\begin{knowledgebox}{读图：框架为什么能提供比通用 IaC 更多的信息}
IaC 知道“创建一个 function”，却未必知道它对应哪个 route、能否静态化、何时 revalidate、是否需要 edge middleware。框架掌握应用语义，平台掌握运行约束；FDI 的优势来自二者在 build time 合并，而不是来自少写几行 YAML。
\end{knowledgebox}

\subsection{Build Output API：把框架适配变成 artifact contract}

Vercel 的 Build Output API 把编译结果定义为文件系统规范：静态资源、functions、routing 和配置被写成可部署 artifact。这样 Next.js、SvelteKit、Nuxt 或其他框架可以生成共同 IR，平台只需消费标准化输出。它也允许本地 build 后使用 prebuilt deployment，从而把 build 与 deploy 解耦。

\lecturefigure{04-build-output-graph.png}{Build Output 把一个应用拆成可部署的 artifact graph。}{Vercel Build Output API 官方文档；课堂 cloud compiler 比喻。}

\begin{importantbox}{编译器抽象带来的四个系统收益}
\begin{enumerate}
  \item \textbf{可重复}：相同输入和版本生成可审计输出；
  \item \textbf{可移植}：框架通过 artifact contract 接入平台；
  \item \textbf{可优化}：平台可按 route 选择 cache、runtime 和 region；
  \item \textbf{可回滚}：部署对象不可变，流量指针可切回旧版本。
\end{enumerate}
\end{importantbox}

\begin{warningbox}{编译器不能推断所有业务意图}
平台可以识别 route 和 cache declaration，却不知道库存数据能否陈旧、某个 API 是否需要强一致、失败是否影响收入。FDI 需要 guard rails 和 override；若把错误默认值自动扩散到全球，自动化只会更快地制造事故。
\end{warningbox}

\subsection{本章小结}

FDI 的本质是从框架语义生成平台 IR，再映射为基础设施。Build Output API 把这条编译边界显式化，使框架生态与平台运行时能够独立演化。

